%==============================================================================================
% General packages before hyperref (no personalized macros here)
%=============================================================================================
\usepackage{amsmath}
	% \usepackage{amsfonts}
	% \usepackage{amssymb} % throws errors with ACM template
	% \usepackage{amsopn} % allows to define new operators, e.g., \argmin
\usepackage{mathtools}
\usepackage{graphicx} % can better scale and rotate graphics than graphic package
\usepackage{url}
\usepackage{enumitem} % remove spaces inside 'itemize' [nolistsep]

% \usepackage{bm} % bold Greek letters
\usepackage{stmaryrd} % for symbol \llceil
	% max 16 math family slots in latex; commented out bm to allow stmaryrd to work

\usepackage{upgreek} % upright Greek letters; together with package "bm" -> upright bold letters
\usepackage{bbding} % for the \CheckmarkBold symbol

\usepackage{listings} % to get sql etc nicely formatted (source code listings)
% \usepackage{fontspec} % to change the SQL font in listings
\lstset{
  basicstyle=\linespread{0.9}\ttfamily,
  columns=fullflexible,
  keepspaces=true,
  mathescape,
  escapechar={|} % https://tex.stackexchange.com/questions/345694/change-color-of-some-items-in-lstlisting  
}



\usepackage{tabularx} % table environment that allows to span a whole line
\usepackage{booktabs} % \toprule
\usepackage{multirow}
\usepackage{hhline} % cannot be overwritten by cellcolor
\usepackage{diagbox}
\usepackage{pgfplotstable} % powerful numerical table macros

	% \usepackage[usenames,table,dvipsnames,x11names]{xcolor} % not needed as options are passed directoy to documentclass 'acmart'
	% \PassOptionsToPackage{table}{xcolor} % xcolor needs to come before tikz which calls xcolor as well. Solution: pass the options before calling tikz

% for colored cells in tables
\usepackage{xcolor,colortbl}    % !!! for some reason, overleaf needs xcolor separately, although that should already have been called by colortbl (10/2019)



\usepackage{tikz} % drawing figures
\usetikzlibrary{matrix}
\usetikzlibrary{calc}
\usetikzlibrary{math}
\usetikzlibrary{arrows.meta,positioning}
\usetikzlibrary{intersections, pgfplots.fillbetween}


\usepackage{easybmat} % \begin{BMAT} alternative syntax for math tabulars and also includes the command addpath which allows to draw a path between matrix elements.

\usepackage{adjustbox}
\def\eox{\unskip\kern 10pt{\unitlength1pt\linethickness{.4pt}$\diamondsuit${}}} % "\eox" command for end of example


\usepackage{rotating}		% for rotated text in tables, \begin{turn}{75}Accuracy\end{turn}

\sloppy % Latex sometimes creates text lines that go into the margin, just to avoid ``ugly'' large spacing. The following prevents this:

\usepackage{xspace} % for macros

% \usepackage[bottom]{footmisc}   % for footnotes





%==============================================================================================
% subfigures
% use package 'subcaption'
% Package 'subfigure' is deprecated and su­per­seded by 'sub­fig' 
% But 'subfig' is not working properly with hyperref.
% 'sub­cap­tion' pack­age is most recent
% See: https://tex.stackexchange.com/questions/13625/subcaption-vs-subfig-best-package-for-referencing-a-subfigure
%=============================================================================================
\usepackage{subcaption} % recommended new package for subfigures (package subfig deprecated)
\captionsetup[figure]{labelfont=bf} % textfont=bf, labelfont=bf, ,font+=bf
\captionsetup[subfigure]{font+=small}
\captionsetup[table]{labelfont=bf} % font+=small, textfont=bf, labelfont=b

	% \captionsetup[subfigure]{justification=centering}
	\captionsetup[figure]{aboveskip=2pt}
	\captionsetup[subfigure]{aboveskip=1pt, belowskip=1pt}		% aboveskip = distance fig and subcaption)
	\captionsetup[figure]{belowskip=-3pt} % space after figures





%==============================================================================================
% Comments and ways to hide them
% macro \note{W: bla bla}: In blue. Start your comment with a short author identifier, e.g. 'W' for Wolfgang. Can be kept in the paper even during and after submission
% macro \todo{W: bla bla}: In red. Those are important issues that need to be fixed. Once they are fixed, replace the \todo with \note
% macro \mkclean: hides the blue comments (not the red ones)
%=============================================================================================
\newcommand{\todo}[1]{{{\color{red}{[\textbf{#1}]}}}} % important comments
\newcommand{\note}[1]{{{\color{magenta}{[#1]}}}} % general comments
\newcommand{\renee}[1]{{{\color{red}{#1}}}}
\newcommand{\rjm}[1]{{{\color{red}{[\textbf{RJM}: #1]}}}}

\newcommand{\revision}[1]{{\color{blue}{#1}}}



% important comments
	% \newcommand{\todo}[1]{{\sffamily{\color{red}{[\small #1]}}}} % needs to be fixed
	% \newcommand{\note}[1]{{\sffamily{\color{blue}{[\small #1]}}}} % comment
% \newcommand{\new}[1]{{\color{violet}{#1}}}

\newcommand{\hide}[1]{} 
\newcommand{\mkclean}{
    \renewcommand{\note}{\hide}
    \renewcommand{\todo}{\hide}
    \renewcommand{\rjm}{\hide}	
}


%\mkclean % comment out this line, to hide all note comments except for those that *need* to be fixed






%==============================================================================================
% ALGORITHM2e
%=============================================================================================
\usepackage[ruled,noend,linesnumbered]{algorithm2e} % boxruled, linesnumbered,ruled

\usepackage{setspace} % for setting line spacing

\DontPrintSemicolon     % does not print ";" at end of each line
	% \LinesNotNumbered
	% \LinesNumbered
\SetNlSty{}{}{}                % style of line numbers "\Setnlsty{small}{}{:}"
\SetAlgoInsideSkip{smallskip}   % distance between caption rule and beginning of algorithm
% FOLLOWING REMOVED FOR VLDB
%\SetAlCapSkip{1.5mm}             % left margin of caption
%\setlength{\algomargin}{1.5mm}    % (1) left margin of algorithm content, but also (!) right margin of caption; also puts line numbers into margin
	% \SetNlSkip{2mm}                % (2) therefore put line numbers back and use "\Indp" twice inside algorithm to put algorithm further to the right
	% \SetAlgoVlined
%\SetInd{1.5mm}{1.5mm}             % corrects for the twice usage of "\Indp"; otherwise vlines are wrong

\SetAlFnt{\small}			% size of actual algorithm text\scriptsize\sf\rm
\SetAlCapFnt{\small}		% size of name of algorithm
\SetAlCapNameFnt{\small}
% ---- For tables ----
\usepackage{array}
\usepackage{makecell}
\usepackage[utf8]{inputenc}
\def\ojoin{\setbox0=\hbox{$\bowtie$}%
  \rule[-.02ex]{.25em}{.4pt}\llap{\rule[\ht0]{.25em}{.4pt}}}
\def\leftouterjoin{\mathbin{\ojoin\mkern-5.8mu\bowtie}}
\def\rightouterjoin{\mathbin{\bowtie\mkern-5.8mu\ojoin}}
\def\fullouterjoin{\mathbin{\ojoin\mkern-5.8mu\bowtie\mkern-5.8mu\ojoin}}

% ---- OPERATORS (requires amsmath) ----



% ---- Algorithm comments ----
%\newcommand{\cc}{olive}
%\newcommand{\algocomment}[1]{\textcolor{\cc}{{//#1}}}
%\providecommand{\cc}{olive}
%\providecommand{\algocomment}[1]{\textcolor{\cc}{{//#1}}}
% ---- Notation ----


% ---- Symbols ----





%==============================================================================================
% CLEVERREF
% better than hyperref's autoref, needs to come after hyperref and algorithm2e (if referencing lines)(Sept 2015)
% \usepackage[nameinlink]{cleveref} 		% option clash with new ACM class (7/2017)
%=============================================================================================
% \usepackage[capitalise,nameinlink]{cleveref} %nameinlink makes whole word colored

% \crefname{algocf}{alg.}{algs.}
% \Crefname{algocf}{Algorithm}{Algorithms}
% \crefalias{AlgoLine}{line}%
% \crefname{proposition}{prop.}{prop.}
% \crefname{figure}{Figure}{Figures}
% \crefname{observation}{Observation}{Observations}

	% \crefformat{equation}{Eq.~(#2#1#3)}						% show reference to equation as (1)
	% \crefrangeformat{equation}{(#3#1#4) to~(#5#2#6)}
	% \crefmultiformat{equation}{(#2#1#3)}%
	% { and~(#2#1#3)}{, (#2#1#3)}{ and~(#2#1#3)}
	% 	% \crefformat{equation}{(#2#1#3)}						% show reference to equation as (1)
	% 	% \crefrangeformat{equation}{(#3#1#4) to~(#5#2#6)}
	% 	% \crefmultiformat{equation}{(#2#1#3)}%
	% 	% { and~(#2#1#3)}{, (#2#1#3)}{ and~(#2#1#3)}
	%
	% % \Crefname{equation}{Equation}{Equations}
	% \Crefformat{equation}{Equation~(#2#1#3)}						% show reference to equation as (1)




	% %==============================================================================
	% % !!! following does not work if \amsthm is used !!!
	% %
	% % These lines fix some annyong property of the ACM template that first lines of definitions
	% % swap over into the margin. Example
	% %   DEFINITION 2.4       (Actual cause).   blabla
	% % instead of
	% %   DEFINITION 2.4  (Actual cause).   blabla
	% %
	% % Solution is to tweak some distance measure in the ACM template. Below separate not to have to tweak the cls file
	% % still has to be tweaked with "~\!\!\!\!" and "~\,\," the ~ is important here (Wolf Dec 2009)
	% % remoeved for new ACM template (7/2017)
	% %==============================================================================
	% \makeatletter
	% \renewcommand{\@opargbegintheorem}[3]{%
	%    \parskip 0pt % GM July 2000 (for tighter spacing)
	%    \trivlist
	%    \item[%
	%        \hskip 10\p@
	%        \hskip \labelsep
	%        {\sc #1\ #2\             %
	%   \setbox\@tempboxa\hbox{(#3)}  %
	%        \ifdim \wd\@tempboxa>\z@ %
	%            \hskip -2\p@\relax    % This number has to be reduced from 5 !!!
	%            \box\@tempboxa       %
	%        \fi.}%
	%    ]
	%    \it
	% 	% \sloppypar	% prevents the overrun into margins of first line for theorem environments for paragraph
	% 	\setlength{\emergencystretch}{6pt}
	% }
	% \makeatother












%==============================================================================================
% New environments (notice that question and answer have slightly different formats)
% NOT NEEDED FOR acmart
% aliascnt still needed for VLDB!
%=============================================================================================
\usepackage{aliascnt}  	
	% ``hyperref’s \autoref command does not work well with theorems that share a counter:
	% it’ll always think it’s a Lemma even if it’s a Remark that shares the Lemma counter.
	% Load this package to fix it. No further intervention needed.''
	% Source: http://absatzen.de/thmtools.html (Jan 2009) 
	% better: http://www.tug.org/applications/hyperref/manual.html (Nov 2009)
	% needs also: thm-patch.sty, parseargs.sty, aliasctr.sty ???
	% see section below the three steps for usage
	% 1 define a new alias counter
	% 2 define the environment (e.g., example)      	
    % 3 reset the counter


% \newtheorem{theorem}{Theorem} % ...[section]

% \newaliascnt{corollary}{theorem}
% \newtheorem{corollary}[corollary]{Corollary}
% \aliascntresetthe{corollary}

% \newaliascnt{example}{theorem}
% \newtheorem{example}[example]{Example}
% \aliascntresetthe{example}

% \newaliascnt{definition}{theorem}
% \newtheorem{definition}[definition]{Definition}
% \aliascntresetthe{definition}

% \newaliascnt{proposition}{theorem}
% \newtheorem{proposition}[proposition]{Proposition}
% \aliascntresetthe{proposition}

% \newaliascnt{lemma}{theorem}
% \newtheorem{lemma}[lemma]{Lemma}
% \aliascntresetthe{lemma}
	
	
% \newaliascnt{conjecture}{theorem}
% \newtheorem{conjecture}[conjecture]{Conjecture}
% \aliascntresetthe{conjecture}

	% \newaliascnt{remark}{theorem}
	% \newtheorem{remark}[remark]{Remark}
	% \aliascntresetthe{remark}
	%
	% \newaliascnt{axiom}{theorem}
	% \newtheorem{axiom}[axiom]{Axiom}
	% \aliascntresetthe{axiom}
	%
	% \newaliascnt{fact}{theorem}
	% \newtheorem{fact}[fact]{Fact}
	% \aliascntresetthe{fact}
	%
	% \newaliascnt{claim}{theorem}
	% \newtheorem{claim}[claim]{Claim}
	% \aliascntresetthe{claim}




%\newtheorem{algo}[theorem]{Algorithm}

\newtheorem{observation}{Observation}
\newtheorem{problem}{Problem}
\newtheorem{answer}{Answer}
\newtheorem{question}{Question}


% ---- For tables ----
\usepackage{array}
\usepackage{makecell}
\usepackage[utf8]{inputenc}
\def\ojoin{\setbox0=\hbox{$\bowtie$}%
  \rule[-.02ex]{.25em}{.4pt}\llap{\rule[\ht0]{.25em}{.4pt}}}
\def\leftouterjoin{\mathbin{\ojoin\mkern-5.8mu\bowtie}}
\def\rightouterjoin{\mathbin{\bowtie\mkern-5.8mu\ojoin}}
\def\fullouterjoin{\mathbin{\ojoin\mkern-5.8mu\bowtie\mkern-5.8mu\ojoin}}

% ---- OPERATORS (requires amsmath) ----



% ---- Algorithm comments ----
\newcommand{\cc}{olive}
\newcommand{\algocomment}[1]{\textcolor{\cc}{{//#1}}}

% ---- Notation ----


% ---- Symbols ----















%==============================================================================================
% Make footnotesizes smaller for VLDB template
%=============================================================================================
% \let\footnotesize\scriptsize
	% \makeatletter
	% \renewcommand\footnotesize{%
	%    \@setfontsize\footnotesize\@ixpt{11}%
	%    \abovedisplayskip 8\p@ \@plus2\p@ \@minus4\p@
	%    \abovedisplayshortskip \z@ \@plus\p@
	%    \belowdisplayshortskip 4\p@ \@plus2\p@ \@minus2\p@
	%    \def\@listi{\leftmargin\leftmargini
	%                \topsep 4\p@ \@plus2\p@ \@minus2\p@
	%                \parsep 2\p@ \@plus\p@ \@minus\p@
	%                \itemsep \parsep}%
	%    \belowdisplayskip \abovedisplayskip
	% }
	% \makeatother









%==============================================================================================
% Change the color use if 'screen=true' (used in acmart.cls around line 482 (4/2019)
%==============================================================================================

\AtEndPreamble{%
    \hypersetup{colorlinks,
      linkcolor=purple,
      citecolor=blue,
      urlcolor=ACMDarkBlue,
      filecolor=ACMDarkBlue}}



%==============================================================================
% These lines fix some annyong property of the old ACM template (= VLDB template) that first lines of definitions
% swap over into the margin. Example
%   DEFINITION 2.4       (Actual cause).   blabla
% instead of 
%   DEFINITION 2.4  (Actual cause).   blabla
% Solution is to tweak some distance measure in the ACM template. Below separate not to have to tweak the cls file
% still has to be tweaked with "~\!\!\!\!" and "~\,\," the ~ is important here (Wolf Dec 2009)
%
% New fix to remove the indentation for theorems environments (Wolf July 2014)
%==============================================================================
%	\makeatletter 
%
%	\renewcommand{\@opargbegintheorem}[3]{%
%	    \parskip 0pt % GM July 2000 (for tighter spacing): 0 pt, 5 pt
%	    \trivlist
%	    \item[%
%	    	\hskip 10\p@										% this number was change to 0, former 10
%	        \hskip \labelsep
%	        {\sc #1\ #2\             % 
%	   \setbox\@tempboxa\hbox{(#3)}  % 
%	        \ifdim \wd\@tempboxa>\z@ % 
%	            \hskip 0\p@\relax    % This number has to be reduced from 5 !!!
%	            \box\@tempboxa       % 
%	        \fi.}%
%	    ]
%	    \it
%	}
%
%	\def\@begintheorem#1#2{%
%	    \parskip 0pt % GM July 2000 (for tighter spacing)
%	    \trivlist
%	    \item[%
%	    	\hskip 10\p@										% this number was change to 0, former 10
%	        \hskip \labelsep
%	        {{\sc #1}\hskip 5\p@\relax#2.}%
%	    ]
%	    \it
%	}
%
%	\makeatother






%=============================================================================================
% Nice backgroundcolor around environments and results
%=============================================================================================
\usepackage{tcolorbox}		% also allows to have shaded backgrounds around environments
\tcbuselibrary{breakable,skins}		% allow to break boxes, allow to use "enhanced jigsaw"


% === Style used for blue example environment
\tcbset{examplestyle/.style={
		enhanced jigsaw,	% if box is broken, don't show border at continuation
		colback=blue!08,	% !20	gray!10
		colframe=blue!08,	% !40
		arc=2mm,
		boxrule=0pt,		% 0pt
		% colframe=red!75!black,
		left=1mm,
		right=1mm,
		left skip= 0mm,  % breaks ACM template unless specified for specific style (4/2019)
		right skip= 0mm, % breaks ACM template unless specified for specific style (4/2019)
		top=1mm,		% 0mm, -1mm
		bottom=1mm,		% problem in package with eq at end of env. just add: \tcbset{bottom=2mm} \tcbset{bottom=0mm}
		breakable,		% allows page break
		parbox = false,		% restores normal text behavior. ELSE: paragraphs are formatted slightly different as the main text. DOWNSIDE: unwanted side effects
		% nobeforeafter,
		% before={\par\pagebreak[0]\smallskip\parindent=0pt},
		before={\par\pagebreak[0]\vspace{1mm}\parindent=0pt},		
		% before={\vspace{1mm}},			
		after={\par\pagebreak[0]\vspace{1mm}\parindent=0pt},				
		% after=5mm,		
		bottomrule = 0mm,
		boxsep = 0mm,					% common padding of hlengthi between the text content and the frame of the box, added to left, right, top, ...
		topsep at break=0pt,			% Additional vertical space at the top of middle and last parts in a break sequence
		bottomsep at break=0pt,			% Additional vertical space at the first and middle parts in a break sequence		
		pad at break=0mm,
		pad before break=1mm,		
		pad after break=1mm,		
		bottomrule at break=0mm,
		toprule at break=0mm,		
		}}

		

\usepackage{footnote}

\tcolorboxenvironment{example}{examplestyle}
%\usepackage{footnote}% allows to wrap an example in "\begin{savenotes}" to force footnotes to the bottom of page


	% \tcbset{
	% 		enhanced jigsaw,	% if box is broken, don't show border at continuation
	% 		% colback=gray!20,
	% 		% colframe=gray!40,
	% 		arc=0mm,
	% 		boxrule=1pt,		% 0pt
	% 		% colframe=red!75!black,
	% 		left=1pt,
	% 		right=1pt,
	% 		% left skip=-1pt,  % breaks ACM template (7/2018)
	% 		% right skip=-1pt, % breaks ACM template (7/2018)
	% 		topsep at break=1mm,			% Additional vertical space at the top of middle and last parts in a break sequence
	% 		top=1pt,		% 0mm, -1mm
	% 		bottom=0mm,		% problem in package with eq at end of env. just add: \tcbset{bottom=2mm} \tcbset{bottom=0mm}
	% 		breakable,		% allows page break
	% 		parbox = false		% restores normal text behavior. ELSE: paragraphs are formatted slightly different as the main text. DOWNSIDE: unwanted side effects
	% 		}

	% \tcolorboxenvironment{example}{}



% === Style used for red result box
\newcommand{\resultbox}[1]{
\begin{tcolorbox}[
	enhanced jigsaw,		% if box is broken, don't show border at continuation
	colback=red!5,
	colframe=red!75!black,	
	arc=0mm,
		% boxrule=1pt,
		% boxrule=0pt,
		% colframe=red!75!black,
	left skip=-1mm,
	right skip=-1mm,	
	left=0mm,
	topsep at break=1mm,			% Additional vertical space at the top of middle and last parts in a break sequence
		% enlarge top by=0mm,				% before box
		% enlarge top at break by=5mm,	% above box, after pagebreak
	right=0mm,
	top=0mm,
	bottom=0mm,		% problem in package with eq at end of env. justadd:\tcbset{bottom=2mm}\tcbset{bottom=0mm}
		% after skip=2mm,
	breakable,		% allows page break
		% before upper={\parindent15pt}
	parbox = false		% restores normal text behavior. ELSE: paragraphs are formatted slightly different as the main text. DOWNSIDE: unwanted side effects
]
\emph{#1}
\end{tcolorbox}
}





%==============================================================================
% defines upright \ell character as \uell (10/2014)
% Source: http://tex.stackexchange.com/questions/175236/typeset-an-upright-ell
%==============================================================================
\makeatletter
\DeclareRobustCommand*\uell{\mathpalette\@uell\relax}
\newcommand*\@uell[2]{
  % We need to adjust the width of \uell to be the same as \ell
  \setbox0=\hbox{$#1\ell$}
  \setbox1=\hbox{\rotatebox{10}{$#1\ell$}}
  \dimen0=\wd0 \advance\dimen0 by -\wd1 \divide\dimen0 by 2
  \mathord{\lower 0.1ex \hbox{\kern\dimen0\unhbox1\kern\dimen0}}
}
\makeatother






%==============================================================================
% allow figures and table to cover 100% of page
% See more: https://aty.sdsu.edu/bibliog/latex/floats.html
%==============================================================================
\renewcommand\dbltopfraction{1} %  max fraction of a two-column page that may be occupied by a (double-wide) float at top
\setcounter{dbltopnumber}{5}%  max nubmers of floats at top of two-column page (dbl = double)









%=============================================================================================
% Rebuttal
% Create marging comments for rebuttal 
% needs \usepackage{soul} for highlighting (12/2018)
%    CHANGE: now replaced soul with ulem
% "\marginpar" does not work for vldb.cls template; requires redefining it with marginnote package "\marginnote" (5/2019)
%=============================================================================================
	% \usepackage{soul} % allows strike-through and highlighting, but very fragile

    % \usepackage[normalem]{ulem} % use normalem to protect \emph; \hl command is more robust
    % \newcommand\hl{\bgroup\markoverwith
    %   {\textcolor{yellow}{\rule[-.5ex]{2pt}{2.5ex}}}\ULon}

% \newcommand{\rebuttal}[2]{\marginpar{\textbf{\textcolor{black}{#1}}}\hl{#2}}
	% \usepackage{marginnote}
	% \newcommand{\rebuttalx}[2]{\marginnote{\textbf{#1}}\hl{#2}}
	% \newcommand{\rebuttalx}[2]{\marginpar{\textbf{#2}}\hl{#1}}	# for some reason does not work: puts text in margin into yellow
	% \newcommand{\rebuttalx}[2]{\marginpar{\textcolor{black}{\textbf{#2}}}\hl{#1}}	# does not work either: puts the comment one line above
%% FOR VLDB.CLS (5/2019) 
\marginparwidth 25pt             	% Margin pars set to 0pt by vldb template
\usepackage{marginnote}
\newcommand{\rebuttal}[2]{\marginnote{\textbf{\textcolor{blue}{#1}}}{\color{blue}#2}}
    % \newcommand{\rebuttal}[2]{\marginnote{\textbf{\textcolor{black}{#1}}}\hl{#2}}

\newcommand{\reviewercomment}[1]{
\begin{tcolorbox}[
	enhanced jigsaw,		% if box is broken, don't show border at continuation
	colback=white!5,
	% colback=blue!20,
	colframe=white!75!black,	
	% colframe=gray!40,
	arc=0mm,
	% boxrule=1pt,
	% boxrule=0pt,
	% colframe=red!75!black,
	left skip=-1mm,
	right skip=1mm,	
	left=0mm,
	topsep at break=1mm,			% Additional vertical space at the top of middle and last parts in a break sequence
		% enlarge top by=0mm,				% before box
		% enlarge top at break by=5mm,	% above box, after pagebreak
	right=0mm,
	top=0mm,
	bottom=0mm,		% problem in package with eq at end of env. justadd:\tcbset{bottom=2mm}\tcbset{bottom=0mm}
		% after skip=2mm,
	breakable,		% allows page break
		% before upper={\parindent15pt}
	parbox = false		% restores normal text behavior. ELSE: paragraphs are formatted slightly different as the main text. DOWNSIDE: unwanted side effects
]
\emph{#1}
\end{tcolorbox}
}

\newcommand{\high}[1]{\textup{\textbf{#1.}}}







%==============================================================================================
% Other macros (add here your personalized macros)
%=============================================================================================
\newcommand{\smallsection}[1]{\vspace{2mm}\noindent\textbf{#1.}} % new subsection type: noindent, bold
\newcommand{\introparagraph}[1]{\textbf{#1.}} % define own new subsection type: noindent, bold (textsc)

\newcommand{\ourmethod}{\textsc{OurMethod}\xspace}
\renewcommand{\epsilon}{\varepsilon} % nicer epsilon symbol
\newcommand{\set}[1]{\{#1\}}                    % Set (as in \set{1,2,3}).
\newcommand{\setof}[2]{\{{#1}\mid{#2}\}}        % Set (as in \setof{x}{x>0}).

\newcommand{\define}{:=} % define symbol
	% \newcommand{\define}{\stackrel{\textup{\tiny def}}{=}}
	% \newcommand{\define}{\triangleq} % define symbol

\newcommand{\datarule}{{\,:\!\!-\,}} % datalog command (March 2010)
\renewcommand{\vec}[1]{\boldsymbol{\mathbf{#1}}}
%\newcommand{\transpose}{\intercal}	% alternatives: \top \intercal

\newcommand{\N}{\mathbb{N}} % Natural numbers
\newcommand{\R}{{\mathbb{R}}} % Real numbers
\newcommand{\bigO}{{\mathcal{O}}} % Big O notation
\newcommand{\bigTH}{{\mathcal{\Theta}}} % Big O notation
\renewcommand{\O}{{\mathcal{O}}} % Big O notation DEPRECATED
\renewcommand{\P}{{\mathbb{P}}} % Probability \P
\newcommand{\PP}[1]{\mathbb{P} [ #1 ]} % Probability \PP{} -> P[], hollow
\newcommand{\E}{\mathbb{E}}
\newcommand{\sql}[1]{\textup{\textsf{\small#1}}}
\newcommand{\datastyle}[1]{\texttt{#1}}


% \newcommand{\val}{\textup{\textrm{cost}}}
\newcommand{\val}{\textup{\textrm{val}}}
% \newcommand{\val}{w}
\newcommand{\var}[1]{{\textup{\texttt{var}}} ( #1 )}  % May 2011 (note naming conflict with variables !!!)
\newcommand{\Var}{\textup{\texttt{Var}}}		






%%%%%%%%%%%%%%%%%%%%%%%%%%%%%%%%%%%%%%%%%%%%%%%%%%%%%%%%%%%%%%%%%%%%%%%%%%%%%%%%
%%%% ORCID icons for authors. Needs to have the image file too in dunnemacros/
%%%%%%%%%%%%%%%%%%%%%%%%%%%%%%%%%%%%%%%%%%%%%%%%%%%%%%%%%%%%%%%%%%%%%%%%%%%%%%%%

% % Use to create orcid icons and links for authors
% % Usage: \orcidicon{<NAME>}{<ORCID>}
% \newcommand{\orcidicon}[2]{%
%   %\hypersetup{urlcolor=black}}{}% don't add colored links to author names
% %  \begingroup\normalfont
%   \href{#2}{#1 \includegraphics[height=\fontcharht\font`\B]{figs/orcid_16x16.png}}%
% }

\usepackage{scalerel}
\usepackage{tikz}
\usetikzlibrary{svg.path}

\definecolor{orcidlogocol}{HTML}{A6CE39}
\tikzset{
  orcidlogo/.pic={
    \fill[orcidlogocol] svg{M256,128c0,70.7-57.3,128-128,128C57.3,256,0,198.7,0,128C0,57.3,57.3,0,128,0C198.7,0,256,57.3,256,128z};
    \fill[white] svg{M86.3,186.2H70.9V79.1h15.4v48.4V186.2z}
                 svg{M108.9,79.1h41.6c39.6,0,57,28.3,57,53.6c0,27.5-21.5,53.6-56.8,53.6h-41.8V79.1z M124.3,172.4h24.5c34.9,0,42.9-26.5,42.9-39.7c0-21.5-13.7-39.7-43.7-39.7h-23.7V172.4z}
                 svg{M88.7,56.8c0,5.5-4.5,10.1-10.1,10.1c-5.6,0-10.1-4.6-10.1-10.1c0-5.6,4.5-10.1,10.1-10.1C84.2,46.7,88.7,51.3,88.7,56.8z};
  }
}

\RequirePackage{etoolbox}
\DeclareRobustCommand\orcidicon[1]{\href{https://orcid.org/#1}{\mbox{\scalerel*{
\begin{tikzpicture}[yscale=-1,transform shape]
\pic{orcidlogo};
\end{tikzpicture}
}{|}}}}

%rjm outerjoins
\def\ojoin{\setbox0=\hbox{$\bowtie$}%
  \rule[-.02ex]{.25em}{.4pt}\llap{\rule[\ht0]{.25em}{.4pt}}}
\def\leftouterjoin{\mathbin{\ojoin\mkern-5.8mu\bowtie}}
\def\rightouterjoin{\mathbin{\bowtie\mkern-5.8mu\ojoin}}
\def\fullouterjoin{\mathbin{\ojoin\mkern-5.8mu\bowtie\mkern-5.8mu\ojoin}}

%rjm nulls
%\def\missing{\mathbin{\perp \perp}}
%\def\missing{\mathbin{\hat{\perp}}}
\def\missing{\mathbin{\pm}}
\def\produced{\mathbin{\perp}}
\def\schema{\mathbin{\mathcal{A}}}
%\def\schema{\mathbin{sc}}



% ---- For tables ----
\usepackage{array}
\usepackage{makecell}
\usepackage[utf8]{inputenc}
\def\ojoin{\setbox0=\hbox{$\bowtie$}%
  \rule[-.02ex]{.25em}{.4pt}\llap{\rule[\ht0]{.25em}{.4pt}}}
\def\leftouterjoin{\mathbin{\ojoin\mkern-5.8mu\bowtie}}
\def\rightouterjoin{\mathbin{\bowtie\mkern-5.8mu\ojoin}}
\def\fullouterjoin{\mathbin{\ojoin\mkern-5.8mu\bowtie\mkern-5.8mu\ojoin}}

% ---- OPERATORS (requires amsmath) ----



% ---- Algorithm comments ----


% ---- Notation ----

%rjm feel free to change the definition of this macro
%\newcommand{\arx}{Archive}
\newcommand{\arx}{arXiv}
\newcommand{\so}{S2ORC}
%\newcommand{\so}{Semantic Scholar}


\newcommand{\ours}{$MLB$}
\newcommand{\transpose}{$MLB_T$}
\newcommand{\header}{$MLB_H$}

% ---- Symbols ----

